% !TEX program = xelatex
% ai-pass-begin: commissioned draft in Jeffrey's writing voice, for review.
\documentclass[12pt,letterpaper]{article}
\usepackage[top=0.8in,bottom=0.8in,left=1.15in,right=1.15in]{geometry}
\usepackage{fontspec}
\setmainfont{Latin Modern Roman}[Extension=.otf,UprightFont=lmroman10-regular,BoldFont=lmroman10-bold,ItalicFont=lmroman10-italic,BoldItalicFont=lmroman10-bolditalic]
\setsansfont{Latin Modern Sans}[Extension=.otf,UprightFont=lmsans10-regular,BoldFont=lmsans10-bold]
\setmonofont{Latin Modern Mono}[Extension=.otf,UprightFont=lmmono10-regular,ItalicFont=lmmono10-italic,Scale=0.85]
\usepackage{xcolor}
\usepackage{titlesec}
\usepackage{enumitem}
\usepackage{fancyhdr}
\usepackage{hyperref}
\usepackage{xurl}
\usepackage{microtype}
\usepackage{ac-paper-essay}
\renewcommand{\essaytitle}[1]{\begin{center}{\acbold\fontsize{28pt}{32pt}\selectfont\color{acdark}#1}\end{center}\vspace{0.3em}}
\linespread{1.20}
\newcommand{\ac}{Aesthetic\textcolor{acpink}{.}Computer}
\hypersetup{colorlinks=true,linkcolor=acpurple,urlcolor=acpurple,pdfauthor={@jeffrey (draft prepared with Codex)},pdftitle={}}
\pagestyle{fancy}
\fancyhf{}
\renewcommand{\headrulewidth}{0pt}
\fancyfoot[C]{\footnotesize\color{acgray}\thepage}
\setlength{\footskip}{24pt}
\raggedbottom
\widowpenalty=10000
\clubpenalty=10000
\begin{document}
\noindent
Rusty, \href{https://rustysantos.substack.com/p/exploring-the-shadows-of-electronic}{you write}: “The dream was recorded but only captured part of the way.” I wonder how much of that missing part is the exchange between a person and the equipment. While I'm making something, every response changes what I do next. The recording keeps a result. It doesn't necessarily tell you what it felt like to get there.

I made a small keyboard instrument called notepat. Press C and it plays C. I want the part your fingers learn to survive a change of computer. The \href{https://papers.aesthetic.computer/notepat-26-arxiv.pdf}{notepat paper} follows the instrument's development; in \href{https://papers.aesthetic.computer/keymaps-social-software-26-arxiv.pdf}{\textit{Keymaps as Social Software}}, I look at the layout itself as a form that travels.\enote{\textit{notepat.com: From Keyboard Toy to System Front Door} traces the instrument; \textit{Keymaps as Social Software}, section 6, describes its shared note layout.}

The layout has its own texture. Notes become reaches, clusters, and repeated shapes under the hand. Some phrases fall easily; others require a stretch. These gestural shapes translate across the technical apparatus. Their sound and timing change, but enough remains for a phrase to be recognized and played again. Meaning and energy travel in the gesture, carried by practice as well as by a file.

That form can be exploited for its familiarity, recalled after an absence, subordinated to another purpose, or brought forward again. A shape I once played as a melody can become accompaniment. Later I can bring it back and hear its earlier role inside the new one. Something of the old gesture's force can return without reproducing the old performance.

But the agreement between versions leaves something out. It says which note I get, not exactly when it arrives or how it stops. A computer coordinates sound, screen updates, and incoming gestures. Latency is the delay before the answer; jitter is how much that delay varies. These are measurable parts of the exchange. They don't explain the whole feeling.\enote{McPherson, Jack, and Moro, \href{https://www.nime.org/proceedings/2016/nime2016_paper0005.pdf}{\textit{Action-Sound Latency: Are Our Tools Fast Enough?}}, NIME 2016.}

I care about a prompt answer because I want to choose the pace. A note can take several seconds to swell and still respond exactly as I intended. An unexpected stall leaves me wondering whether the computer received my gesture. I can reduce an unwanted delay and still make slow music.

Drawing works this way too. I put down a mark, look at it, and answer it with another. If it appears late, I may have moved my hand before I see what I've done. The delay can change the next decision without becoming something a viewer can identify in the finished image.\enote{On drawing as a conversation with a recording: my \href{https://www.youtube.com/watch?v=K8xCBPW1rmo}{lecture on mediation}, April 14, 2020.}

\clearpage
\noindent
You also write: “An artist’s role is to elevate the content above the container.” My work pulls me toward the container, but also toward the forms that pass between containers. The layout gives a gesture somewhere to return. Each instrument gives it particular resistance, timing, and sound. I work on those conditions so somebody else can make music I couldn't predict.

This is why I maintain several implementations of my instrument across different software and hardware conditions. Some parts are shared, but each has different code fitted to its machine's material environment. I want to use the apparatus's qualities natively. The shared layout lets a player carry something learned into those different conditions.

The website, \href{https://notepat.com}{notepat.com}, was the first implementation. It is made for mobile phones and quick, casual play: open a link and touch the notes. Its availability comes with the browser's constraints. Of these versions, it is where I accept the most delay between gesture and sound; the phone and browser determine how much. It puts the instrument within easy reach.

\href{https://papers.aesthetic.computer/ac-native-os-26-arxiv.pdf}{ACOS}—the \ac{} operating system—takes another approach. On supported older, inexpensive hardware, it boots straight into notepat and pursues very low latency. An old laptop can become a dedicated instrument. The code works close to the hardware so the machine can give its attention to playing.\enote{For the hardware-dependent latency work, see \href{https://papers.aesthetic.computer/where-the-microseconds-go-26-arxiv.pdf}{\textit{Where the Microseconds Go}}.}

Menu Band, at \href{https://menuband.app}{menuband.app}, lives casually in the Mac's menu bar. It uses the Mac's General MIDI instrument palette, audio interfaces, and native audio and MIDI routing. I can play a phrase, record it, export audio and MIDI, or send notes into recording applications such as \href{https://help.ableton.com/hc/en-us/articles/211476789-Setting-up-an-Audio-Interface}{Ableton Live}. The instrument joins the musical environment already on that computer.

So vendor systems and closed-source operating systems belong on the playing field too. A particular interface, driver, sound library, or recording application may be part of why someone chooses that apparatus. They all need to play. The material conditions are chosen by the musicker. I maintain different code to meet those conditions, while keeping the layout available to translate between them.

% ai-pass-end
\end{document}
